% Part: first-order-logic
% Chapter: beyond
% Section: other-logics

\documentclass[../../../include/open-logic-section]{subfiles}

\begin{document}

\olfileid{fol}{byd}{oth}

\olsection{ہور منطق دیاں قسماں}

ہُن تک شاید تہاڈے ذہن وچ ایہ گل آ گئی ہووے کہ نویں منطق بنانی اوکھی نہیں۔
تُسی وی اپنی نحوی ساخت بنا سکدے او، اک استدلالی نظام گھڑ سکدے او، تے اوہدے
نال ملدی معنیات ترتیب دے سکدے او۔ جے تسی چاہندے او کہ !!{derivation}
دا نظام انہاں معنیات دے لحاظ نال تمام ہوئے، تاں کجھ سوجھ بوجھ درکار ہو
سکدی اے؛ تے عام لوکاں نوں ایہ منوان لئی وی محنت لگ سکدی اے کہ تہاڈی
منطق واقعی دلچسپ اے۔ پر ایس دے بدلے تُسی اپنی منطق دیاں ریاضیاتی تے
کمپیوٹیشنل خاصیتاں دی کھوج کردیاں کئی گھنٹے چنگی، بے ضرر تفریح وچ
گزار سکدے او۔

پچھلیاں کئی دہاکیاں وچ رسمی منطق دیاں قسماں حیران کن حد تک ودھیاں نیں۔
دھندلی منطق (fuzzy logic) مبہم خاصیتاں بارے استدلال دا نمونہ بناؤن لئی
بنائی جاندی اے۔ احتمالاتی منطق (probabilistic logic) غیر یقینی صورتاں
بارے استدلال دا نمونہ بناندی اے۔ ڈیفالٹ منطق دے نظام (default logics)
تے نویں معلومات نال نتیجے پلٹ سکدے ہون والی منطق دے نظام
(nonmonotonic logics) ایسے استدلال دا نمونہ بناندے نیں: یعنی اوہ ”معقول“ نتیجے جو نویں
معلومات آون اُتے رد ہو سکدے نیں۔ علم بارے استدلال لئی علمی منطق دیاں
قسماں (epistemic logics)، سبب تے نتیجے دے رشتیاں بارے استدلال لئی
سببیتی منطق دیاں قسماں (causal logics)، تے اخلاقی فرضاں تے ذمہ واریاں
بارے استدلال لئی ”اخلاقی فرض دی منطق“ (deontic logics) وی نیں۔ ایہناں
نظاماں نوں متعارف کراؤن دی اصل ترغیب فلسفیانہ، ریاضیاتی یا
کمپیوٹیشنل ہووے، ایس دے مطابق تہاڈے سامنے انہاں دا مطالعہ ریاضیاتی
منطق، فلسفیانہ منطق، مصنوعی ذہانت، ادراکی سائنس یا کسے ہور شعبے دے
تحت آ سکدا اے۔

ایہ فہرست لمی ہوندی جاندی اے، تے امکان بے انت لگدے نیں۔ سارے انسانی
استدلال نوں حسابی عمل وچ ڈھال دیون دا لائبنز دا خواب شاید اسیں کدے
پورا نہ کر سکیے---پر ایہ گل سانوں کوشش کرن توں نہیں روک سکدی۔

\end{document}
